\documentclass[a4paper,12pt]{article}

\usepackage[T2A]{fontenc}
\usepackage[utf8]{inputenc}
\usepackage[russian]{babel}

\usepackage{amsmath,amssymb,mathtools}
\usepackage{helvet}
\renewcommand{\familydefault}{\sfdefault}
\usepackage{icomma}

\usepackage[a4paper,margin=2.1cm,headheight=15pt]{geometry}
\usepackage{microtype}
\usepackage{tikz}
\usepackage{tkz-euclide}
\usepackage{tabularx,booktabs}
\usepackage{enumitem}
\usepackage{hyperref}
\usepackage{parskip}
\usepackage{fancyhdr}
\usepackage{setspace}
\usepackage{xcolor}

% ---------- Палитра ----------
\definecolor{accent}{HTML}{008080}
\definecolor{darkaccent}{HTML}{004D4D}
\definecolor{textgray}{HTML}{333333}
\definecolor{softgray}{HTML}{EEF3F3}
\definecolor{linegray}{HTML}{B9C8C8}

% ---------- Общая типографика ----------
\onehalfspacing
\setlength{\parindent}{0pt}
\setlength{\parskip}{0.55em}
\setlist[itemize]{leftmargin=1.5em,itemsep=0.25em,topsep=0.3em}
\setlist[enumerate]{leftmargin=1.8em,itemsep=0.55em,topsep=0.4em}
\renewcommand{\arraystretch}{1.28}

\hypersetup{
    colorlinks=true,
    linkcolor=darkaccent,
    urlcolor=darkaccent
}

\pagestyle{fancy}
\fancyhf{}
\fancyhead[L]{\small\color{textgray}Векторы и скалярное произведение}
\fancyhead[R]{\small\color{textgray}Николь Саркисьянц}
\fancyfoot[C]{\small\color{textgray}\thepage}
\renewcommand{\headrulewidth}{0.4pt}

\newcommand{\vect}[1]{\overrightarrow{#1}}
\newcommand{\vct}[1]{\vec{#1}}
\newcommand{\sectionline}{%
    \par\noindent
    {\color{linegray}\rule{\linewidth}{0.5pt}}
    \par
}

\begin{document}

% =========================================================
% ТИТУЛЬНЫЙ ЛИСТ
% =========================================================

\begin{titlepage}
\thispagestyle{empty}

\begin{center}

\vspace*{2.8cm}

{\Large\color{darkaccent}\textbf{ЧЕК-ЛИСТ}}

\vspace{0.8cm}

{\fontsize{25}{30}\selectfont
\color{accent}\textbf{Векторы и скалярное произведение}}

\vspace{0.35cm}

{\Large\color{textgray}
Координатный метод в геометрии}

\vspace{1.6cm}

{\large
Ключевые формулы, алгоритмы и приёмы\\
для задач ЕГЭ и задач повышенной сложности}

\vfill

{\large
\textbf{Лёвин Артём Александрович}\\[0.2cm]
эксклюзивно для Николь Саркисьянц}

\vspace{0.8cm}

{\color{textgray}\today}

\vspace{1.8cm}

\end{center}

% Обязательная титульная кривая Безье
\begin{tikzpicture}[remember picture,overlay]
    \draw[
        color=accent,
        line width=1.2pt
    ]
    ([xshift=1.9cm,yshift=1.6cm]current page.south west)
    .. controls
    ([xshift=6.2cm,yshift=2.5cm]current page.south west)
    and
    ([xshift=-6.5cm,yshift=0.7cm]current page.south east)
    ..
    ([xshift=-1.9cm,yshift=1.7cm]current page.south east);
\end{tikzpicture}

\end{titlepage}

% =========================================================
% ОСНОВНОЙ ЧЕК-ЛИСТ
% =========================================================

\section*{\color{darkaccent}1. Базовые действия с векторами}

Пусть
\[
\vct a=(x_a;y_a),
\qquad
\vct b=(x_b;y_b).
\]

Тогда:

\[
\vct a+\vct b=(x_a+x_b;\,y_a+y_b),
\]

\[
\vct a-\vct b=(x_a-x_b;\,y_a-y_b),
\]

\[
k\vct a=(kx_a;\,ky_a).
\]

\textbf{Главный принцип:} все линейные действия с векторами выполняются
покоординатно.

\subsection*{\color{accent}Вектор по двум точкам}

Если
\[
A(x_A;y_A),
\qquad
B(x_B;y_B),
\]
то
\[
\boxed{\vect{AB}=(x_B-x_A;\,y_B-y_A)}.
\]

То есть всегда используйте порядок

\[
\boxed{\text{конец} - \text{начало}}.
\]

Для пространства аналогично:
\[
\vect{AB}
=
(x_B-x_A;\,y_B-y_A;\,z_B-z_A).
\]

\sectionline

\section*{\color{darkaccent}2. Длина вектора}

Для
\[
\vct a=(x;y)
\]
имеем
\[
\boxed{|\vct a|=\sqrt{x^2+y^2}}.
\]

В пространстве:
\[
\boxed{|\vct a|=\sqrt{x^2+y^2+z^2}}.
\]

Если координата отрицательная, при возведении в квадрат удобно
явно использовать скобки:
\[
(-7)^2=49.
\]

\subsection*{\color{accent}Пример}

Пусть
\[
\vct a=(-3;2),
\qquad
\vct b=(-1;9).
\]

Найдём
\[
\vct c=3\vct a-2\vct b.
\]

Последовательно:
\[
3\vct a=(-9;6),
\qquad
2\vct b=(-2;18),
\]
поэтому
\[
\vct c=(-9;6)-(-2;18)=(-7;-12).
\]

Следовательно,
\[
|\vct c|
=
\sqrt{(-7)^2+(-12)^2}
=
\boxed{\sqrt{193}}.
\]

\textbf{Экзаменационный совет:}
не выполняйте несколько знаковых операций мысленно в одну строку.
Запись промежуточных координат резко уменьшает риск арифметической ошибки.

\sectionline

\section*{\color{darkaccent}3. Скалярное произведение}

Скалярное произведение можно вычислять двумя основными способами.

\subsection*{\color{accent}Через координаты}

Для
\[
\vct a=(x_a;y_a),
\qquad
\vct b=(x_b;y_b)
\]
получаем
\[
\boxed{
\vct a\cdot\vct b=x_ax_b+y_ay_b
}.
\]

В пространстве:
\[
\boxed{
\vct a\cdot\vct b=x_ax_b+y_ay_b+z_az_b
}.
\]

\subsection*{\color{accent}Через длины и угол}

Для ненулевых векторов
\[
\boxed{
\vct a\cdot\vct b
=
|\vct a|\,|\vct b|\cos\varphi
},
\]
где
\[
0^\circ\le\varphi\le180^\circ
\]
--- угол между векторами.

\begin{table}[h]
\centering
\caption{Как выбирать формулу скалярного произведения}
\begin{tabularx}{\linewidth}{@{}lX@{}}
\toprule
\textbf{Что известно} & \textbf{Обычно удобнее использовать} \\
\midrule
Координаты обоих векторов &
$x_ax_b+y_ay_b$ \\

Длины и угол между векторами &
$|\vct a|\,|\vct b|\cos\varphi$ \\

Нужно определить угол по координатам &
сначала скалярное произведение, затем формулу для $\cos\varphi$ \\
\bottomrule
\end{tabularx}
\end{table}

\sectionline

\section*{\color{darkaccent}4. Перпендикулярность и знак скалярного произведения}

Для ненулевых векторов:
\[
\boxed{
\vct a\perp\vct b
\iff
\vct a\cdot\vct b=0
}.
\]

Это один из основных способов составления уравнения с параметром.

Например, если
\[
\vct a=(x;3),
\qquad
\vct b=(-1;4),
\]
то условие перпендикулярности даёт
\[
-x+12=0,
\]
откуда
\[
x=12.
\]

Кроме того,

\[
\vct a\cdot\vct b
\begin{cases}
>0, & \varphi<90^\circ,\\
=0, & \varphi=90^\circ,\\
<0, & \varphi>90^\circ.
\end{cases}
\]

\textbf{Важно:}
угол между векторами \emph{не обязан быть острым}.
Если скалярное произведение отрицательно, угол между векторами тупой.

\sectionline

\section*{\color{darkaccent}5. Как найти угол между векторами}

Из формулы скалярного произведения:
\[
\boxed{
\cos\varphi
=
\frac{\vct a\cdot\vct b}
{|\vct a|\,|\vct b|}
}.
\]

После вычисления косинуса выбирается единственный угол
\[
\varphi\in[0^\circ;180^\circ].
\]

Например, если
\[
|\vct a|=2,
\qquad
|\vct b|=5,
\qquad
\vct a\cdot\vct b=-5,
\]
то
\[
\cos\varphi
=
\frac{-5}{2\cdot5}
=
-\frac12.
\]

Следовательно,
\[
\boxed{\varphi=120^\circ}.
\]

Заменять $-\frac12$ на $\frac12$ здесь нельзя:
знак косинуса содержит информацию о том, является угол острым или тупым.

\sectionline

\section*{\color{darkaccent}6. Квадрат длины суммы и разности}

Очень важные формулы:
\[
\boxed{
|\vct a-\vct b|^2
=
|\vct a|^2
+
|\vct b|^2
-
2\vct a\cdot\vct b
}
\]

и
\[
\boxed{
|\vct a+\vct b|^2
=
|\vct a|^2
+
|\vct b|^2
+
2\vct a\cdot\vct b
}.
\]

Если известен угол,
\[
\vct a\cdot\vct b
=
|\vct a|\,|\vct b|\cos\varphi.
\]

\subsection*{\color{accent}Пример}

Пусть
\[
|\vct a|=6,
\qquad
|\vct b|=8,
\qquad
\varphi=60^\circ.
\]

Сначала
\[
\vct a\cdot\vct b
=
6\cdot8\cdot\cos60^\circ
=
48\cdot\frac12
=
24.
\]

Тогда
\[
|\vct a-\vct b|^2
=
36+64-2\cdot24
=
\boxed{52}.
\]

Если требуется именно длина:
\[
|\vct a-\vct b|
=
\sqrt{52}
=
\boxed{2\sqrt{13}}.
\]

\textbf{Не путайте} длину и квадрат длины.

\sectionline

\section*{\color{darkaccent}7. Координатный метод: координаты не обязаны быть даны}

Если в условии отсутствует координатная плоскость, её часто можно
\textbf{ввести самостоятельно}.

Особенно выгодно это делать, если в фигуре есть:

\begin{itemize}
    \item перпендикулярные прямые;
    \item диагонали, пересекающиеся под прямым углом;
    \item середины отрезков;
    \item оси симметрии;
    \item высоты.
\end{itemize}

\subsection*{\color{accent}Базовый алгоритм}

\begin{enumerate}
    \item Выберите начало координат в удобной характерной точке.
    \item Направьте оси вдоль перпендикулярных прямых.
    \item Запишите координаты ключевых точек.
    \item Получите координаты нужных векторов по формуле
    \[
    \vect{AB}=B-A.
    \]
    \item Примените скалярное произведение, длину или формулу косинуса.
\end{enumerate}

\subsection*{\color{accent}Пример: ромб}

Пусть диагонали ромба имеют длины
\[
AC=8,
\qquad
BD=6.
\]

Диагонали ромба взаимно перпендикулярны и делятся точкой пересечения
пополам. Разместим начало координат в точке $O$.

\begin{center}
\begin{tikzpicture}[scale=0.68]
    \tkzSetUpLine[line width=1pt,color=accent]
    \tkzSetUpPoint[color=darkaccent,fill=darkaccent,size=3]
    \tkzSetUpLabel[color=black,font=\small]

    \tkzDefPoint(-4,0){A}
    \tkzDefPoint(0,3){B}
    \tkzDefPoint(4,0){C}
    \tkzDefPoint(0,-3){D}
    \tkzDefPoint(0,0){O}

    \tkzDrawPolygon(A,B,C,D)
    \tkzDrawSegments(A,C B,D)

    \tkzDrawPoints(A,B,C,D,O)

    \tkzLabelPoints[left](A)
    \tkzLabelPoints[right](C)
    \tkzLabelPoints[above](B)
    \tkzLabelPoints[below](D)
    \tkzLabelPoints[below right](O)

    \tkzMarkRightAngle[size=0.28](C,O,B)
\end{tikzpicture}
\end{center}

Получаем
\[
A(-4;0),
\qquad
B(0;3),
\qquad
D(0;-3).
\]

Тогда
\[
\vect{AB}=(4;3),
\qquad
\vect{AD}=(4;-3).
\]

Поэтому
\[
\vect{AB}\cdot\vect{AD}
=
4\cdot4+3\cdot(-3)
=
16-9
=
\boxed{7}.
\]

Кроме того,
\[
|\vect{AB}|=|\vect{AD}|=5,
\]
поэтому для угла $\alpha=\angle BAD$:
\[
\cos\alpha
=
\frac{7}{5\cdot5}
=
\boxed{\frac7{25}}.
\]

\textbf{Смысл приёма:}
координатный метод переводит геометрическую задачу в небольшой набор
стандартных алгебраических операций.

\sectionline

\section*{\color{darkaccent}8. Середина отрезка}

Если
\[
A(x_1;y_1),
\qquad
C(x_2;y_2),
\]
а $M$ --- середина $AC$, то
\[
\boxed{
M\left(
\frac{x_1+x_2}{2};
\frac{y_1+y_2}{2}
\right)
}.
\]

В векторной форме:
\[
\vct{OM}
=
\frac{\vct{OA}+\vct{OC}}{2}.
\]

Это особенно полезно в задачах с медианами.

\sectionline

\section*{\color{darkaccent}9. Полезная формула для медианы}

Пусть в треугольнике $ABC$ точка $M$ --- середина стороны $AC$.

Тогда
\[
\boxed{
\vect{BM}\cdot\vect{AC}
=
\frac{BC^2-AB^2}{2}
}.
\]

Эта формула позволяет решить ряд задач значительно короче, чем
через явный поиск высоты и координат всех точек.

\subsection*{\color{accent}Почему формула работает}

Положим
\[
\vect{AB}=\vct u,
\qquad
\vect{AC}=\vct v.
\]

Так как $M$ --- середина $AC$,
\[
\vect{AM}=\frac12\vct v.
\]

Следовательно,
\[
\vect{BM}
=
\vect{BA}+\vect{AM}
=
-\vct u+\frac12\vct v.
\]

Поэтому
\[
\vect{BM}\cdot\vect{AC}
=
-\vct u\cdot\vct v+\frac12|\vct v|^2.
\]

Из
\[
BC^2
=
|\vct v-\vct u|^2
=
AC^2+AB^2-2\vct u\cdot\vct v
\]
следует
\[
\vct u\cdot\vct v
=
\frac{AB^2+AC^2-BC^2}{2}.
\]

Подстановка даёт
\[
\vect{BM}\cdot\vect{AC}
=
\frac{BC^2-AB^2}{2}.
\]

\subsection*{\color{accent}Пример}

Если
\[
AB=5,\qquad BC=6,\qquad AC=7,
\]
то
\[
\vect{BM}\cdot\vect{AC}
=
\frac{6^2-5^2}{2}
=
\frac{36-25}{2}
=
\boxed{\frac{11}{2}}.
\]

\sectionline

\section*{\color{darkaccent}10. Если всё же нужен классический координатный путь}

При введении высоты удобно использовать стандартную цепочку:

\[
\boxed{
\text{три стороны}
\longrightarrow
\text{площадь по Герону}
\longrightarrow
\text{высота}
\longrightarrow
\text{проекции}
\longrightarrow
\text{координаты}
}.
\]

Если стороны треугольника равны $a$, $b$, $c$, то
\[
p=\frac{a+b+c}{2},
\]
\[
\boxed{
S=\sqrt{p(p-a)(p-b)(p-c)}
}.
\]

Если к стороне $c$ проведена высота $h$, то
\[
S=\frac12ch,
\]
поэтому
\[
\boxed{h=\frac{2S}{c}}.
\]

После этого в образовавшихся прямоугольных треугольниках можно
использовать теорему Пифагора.

\textbf{Практический ориентир:}
если в геометрической задаче Вы не видите следующий шаг, проверьте,
можно ли построить или обнаружить прямоугольный треугольник.

\sectionline

\section*{\color{darkaccent}11. Что особенно важно не перепутать}

\begin{enumerate}
    \item
    Координаты $\vect{AB}$ считаются как
    \[
    B-A,
    \]
    а не наоборот.

    \item
    При умножении вектора на число умножаются \textbf{все} координаты.

    \item
    Для разности векторов знаки обрабатываются покоординатно.

    \item
    В формуле
    \[
    |\vct a-\vct b|^2
    =
    |\vct a|^2+|\vct b|^2-2\vct a\cdot\vct b
    \]
    перед удвоенным скалярным произведением стоит минус.

    \item
    Если требуется длина, а Вы нашли её квадрат, необходимо извлечь
    \textbf{неотрицательный} квадратный корень.

    \item
    Угол между ненулевыми векторами находится в диапазоне
    \[
    0^\circ\le\varphi\le180^\circ.
    \]

    \item
    Отрицательный косинус означает тупой угол, а не необходимость
    отбросить минус.

    \item
    Если
    \[
    \vct a\cdot\vct b=0,
    \]
    то ненулевые векторы перпендикулярны.

    \item
    Отсутствие готовых координат в условии не запрещает координатный
    метод: систему координат можно ввести самостоятельно.

    \item
    В сложной геометрической задаче сначала ищите структуру:
    перпендикулярность, середины, симметрию, высоты и удобные оси.
\end{enumerate}

% =========================================================
% ТРЕНИРОВОЧНЫЙ БЛОК
% =========================================================

\clearpage
\section*{\color{darkaccent}Тренировочный блок}

\textbf{Рекомендация:}
решайте задачи письменно и не пропускайте промежуточные координатные
вычисления.

\subsection*{\color{accent}Средний уровень}

\begin{enumerate}[label=\textbf{\arabic*.}]

\item
Даны векторы
\[
\vct a=(-2;5),
\qquad
\vct b=(3;-1).
\]
Найдите координаты и длину вектора
\[
2\vct a-\vct b.
\]

\item
Известно, что
\[
|\vct a|=4,
\qquad
|\vct b|=7,
\qquad
\angle(\vct a,\vct b)=60^\circ.
\]
Найдите
\[
\vct a\cdot\vct b
\]
и
\[
|\vct a-\vct b|^2.
\]

\item
При каком значении $x$ векторы
\[
\vct a=(x;3),
\qquad
\vct b=(2;-4)
\]
перпендикулярны?

\end{enumerate}

\subsection*{\color{accent}Выше среднего}

\begin{enumerate}[label=\textbf{\arabic*.},resume]

\item
Даны
\[
\vct a=(1;-2),
\qquad
\vct b=(4;3).
\]
Найдите
\[
\cos\varphi,
\]
где $\varphi$ --- угол между векторами, и определите, является ли этот
угол острым, прямым или тупым.

\item
Известно:
\[
|\vct a|=5,
\qquad
|\vct b|=6,
\qquad
\vct a\cdot\vct b=-15.
\]
Найдите угол между векторами.

\item
Диагонали ромба имеют длины $10$ и $24$.
Найдите скалярное произведение двух векторов, соответствующих
смежным сторонам ромба и имеющих общее начало в одной вершине.

\item
В треугольнике $ABC$
\[
AB=7,\qquad BC=9,\qquad AC=10.
\]
Точка $M$ --- середина стороны $AC$.
Найдите
\[
\vect{BM}\cdot\vect{AC}.
\]

\item
Известно:
\[
|\vct a|=3,
\qquad
|\vct b|=4,
\qquad
\angle(\vct a,\vct b)=120^\circ.
\]
Найдите
\[
|2\vct a+\vct b|.
\]

\item
При каком значении $t$ пространственные векторы
\[
\vct a=(t;1;2),
\qquad
\vct b=(2;-1;3)
\]
перпендикулярны?

\end{enumerate}

\subsection*{\color{accent}Сложная задача}

\begin{enumerate}[label=\textbf{\arabic*.},resume]

\item
В треугольнике $ABC$
\[
AB=8,\qquad BC=10,\qquad AC=12.
\]
Точка $M$ --- середина $AC$, а точка $N$ --- середина $AB$.

Найдите
\[
\vect{BM}\cdot\vect{CN}.
\]

\end{enumerate}

% =========================================================
% ОТВЕТЫ
% =========================================================

\clearpage
\section*{\color{darkaccent}Ответы к тренировочному блоку}

\begin{table}[h]
\centering
\caption{Краткие ответы}
\begin{tabularx}{\linewidth}{@{}cX@{}}
\toprule
\textbf{№} & \textbf{Ответ} \\
\midrule

1 &
\[
2\vct a-\vct b=(-7;11),
\qquad
|2\vct a-\vct b|=\sqrt{170}.
\]
\\

2 &
\[
\vct a\cdot\vct b=14,
\qquad
|\vct a-\vct b|^2=37.
\]
\\

3 &
\[
x=6.
\]
\\

4 &
\[
\cos\varphi=-\frac{2}{5\sqrt5}
=
-\frac{2\sqrt5}{25};
\qquad
\varphi\text{ --- тупой}.
\]
\\

5 &
\[
120^\circ.
\]
\\

6 &
\[
-119.
\]
\\

7 &
\[
16.
\]
\\

8 &
\[
2\sqrt7.
\]
\\

9 &
\[
t=-\frac52.
\]
\\

10 &
\[
-\frac{73}{2}.
\]
\\

\bottomrule
\end{tabularx}
\end{table}

\vfill

\begin{center}
{\small\color{textgray}
Проверяйте не только итоговый ответ, но и выбор формулы,
направление вектора и знаки координат.}
\end{center}

\end{document}